% ---------------------------------------------------------------------------
% Introduction (input for theory_section.tex)
% ---------------------------------------------------------------------------
In Ref.~\cite{MelladoMunoz2020} we immersed an evaporating $^{87}$Rb cloud in a
magneto-optical trap of $^{39}$K and observed condensed fractions above the
equilibrium value for the measured atom number and temperature, up to
$F\simeq0.5$ at $\tau=T/T_c\simeq1$, together with a relaxation of the excess
fraction on a time scale of $3.9$~s after the K atoms were removed. The
comparison was made with the equilibrium relation $F(\tau)$ of an ideal gas
with the leading interaction correction, and the dynamics was described by a
two-mode rate model whose growth and loss coefficients, $\bar W$ and $\bar K$,
and effective thermal population were fitted to the data. The quantum kinetic
theory of condensate growth~\cite{Gardiner1997,Gardiner1998} provides a
microscopic expression for these coefficients, but it treats the thermal cloud
as a reservoir of fixed temperature and chemical potential. In the experiment
neither assumption holds: the thermal cloud contains only $10^4$ atoms, loses
half of them in a second, and cools by a factor of three during the window in
which the condensate forms. The ergodic kinetic theory of
Refs.~\cite{Zaremba1999,Bijlsma2000} removes both assumptions. It was applied to
the growth of a condensate after a sudden evaporative cut in a closed trap,
where it reproduced the growth curves of the MIT experiment within a factor of
two to three in the rate, but it has not been applied to a continuously
evaporating gas, to a gas with an energy-independent loss acting on both
components, or to the question of whether a condensate can persist in contact
with an undersaturated thermal cloud.

In this work we solve the ergodic kinetic equations for the complete protocol
of Ref.~\cite{MelladoMunoz2020}: the last second of the evaporation ramp with
and without the K atoms, the ramp with four times the K density, and the hold
after the K atoms are removed. The elastic Rb-Rb collision physics enters with
no free parameter; the evaporation ramp is calibrated once on the data taken
without K, and the data with K are then compared with the model without further
adjustment, except for the strength of the K-induced loss, which we take from
the measured atom numbers. We ask three questions. (i) Does the model
reproduce the atom number, temperature and condensed fraction of the
distillation data? (ii) What does the model predict for the lifetime of the
measured supercritical states? (iii) Which part of the measured deviation from
the equilibrium curve is a consequence of comparing with an equilibrium at
fixed temperature rather than at fixed energy?
